% --- Conteúdo CORRIGIDO para o arquivo Algoritmos/bubble_sort.tex ---
% --- Conteúdo ATUALIZADO para Algoritmos/bubble_sort.tex ---

\qquad O \textit{Bubble Sort} (ordenação por bolha) é um algoritmo de ordenação simples, geralmente utilizado para fins didáticos. Seu funcionamento consiste em percorrer o vetor diversas vezes, comparando pares de elementos adjacentes e trocando-os de posição quando estão fora de ordem. A cada passagem, o maior elemento ainda não ordenado é deslocado para o final do vetor. O algoritmo termina quando uma passagem completa não realiza nenhuma troca, o que indica que o vetor já está ordenado \cite{cormen}.

% ... (resto do arquivo com a tabela de custos) ...

\qquad A seguir é apresentado o código-fonte da implementação em C++, junto com a análise do custo de cada instrução e do número de vezes que ela é executada, considerando um vetor de entrada de tamanho $n$.

\vspace{1cm} % Adiciona um espaço vertical

\begin{center}
\begin{tabular}{l l l l}
\hline
\multicolumn{2}{l}{\textbf{void bubbleSort(vector<int>\& A)}} & \textbf{Custo} & \textbf{Vezes} \\
\hline
1 & \texttt{int n = A.size();} & $c_1$ & 1 \\
2 & \texttt{for (int i = 0; i < n - 1; i++) \{} & $c_2$ & $n$ \\
3 & \texttt{    for (int j = 0; j < n - i - 1; j++) \{} & $c_3$ & $\frac{n(n-1)}{2}$ \\
4 & \texttt{        if (A[j] > A[j + 1]) \{} & $c_4$ & $\frac{n(n-1)}{2}$ \\
5 & \texttt{            int temp = A[j];} & $c_5$ & (depende da entrada) \\
6 & \texttt{            A[j] = A[j + 1];} & $c_6$ & (depende da entrada) \\
7 & \texttt{            A[j + 1] = temp;} & $c_7$ & (depende da entrada) \\
8 & \texttt{        \}} & 0 & - \\
9 & \texttt{    \}} & 0 & - \\
10 & \texttt{\}} & 0 & - \\
\hline
\end{tabular}
\end{center}

\vspace{1cm} % Adiciona um espaço vertical

\qquad Os valores $c_1, c_2, \dots, c_7$ representam o custo de cada linha do algoritmo. Os laços aninhados resultam em uma soma de execuções que se aproxima de $\frac{n(n-1)}{2}$. Já as linhas de troca (5, 6 e 7) só são executadas quando os elementos comparados estão fora de ordem, portanto o número de execuções dessas linhas depende da entrada.